\section{Double Q-Learning y el problema de sobreestimación}

\subsection{Sesgo de sobreestimación}

$$\max_{a'} Q_\phi(s', a') \geq Q^*(s', \pi^*(s'))$$

Tomar la operación max sobreestima sistemáticamente el valor de Q, porque elige la \textbf{dirección con más ruido}.

\subsection{Double DQN}

Se usa la red actual para elegir la acción, pero se evalúa con la red target:
$$y = r + \gamma Q_{\phi'}\left(s', \arg\max_{a'} Q_\phi(s', a')\right)$$

Esto mitiga en cierta medida el problema de sobreestimación.

\subsection{Resumen de la sección}

Q-Learning tiene un sesgo sistemático de sobreestimación. Double Q-Learning mitiga este problema al separar la elección de la acción de su evaluación.

